\begin{topic}[id=tensorflow_intro,
             difficulty={Mittel},
             type={Systemanalyse + Validierung},
             prerequisites={Grundlagen Python; Basisverständnis maschinelles Lernen und Modellinferenz},
             practice={optional},
             owner={Dozententeam},
             updated={2026-04-13},
             tags={ LiteRT, TensorFlow Lite, Edge AI, Konvertierung, Quantisierung, Operatoren, Validierung }]{LiteRT auf Edge-Geräten: Konvertierung, Quantisierung und Validierung}

\TopicDescription{Dieses Thema betrachtet nicht TensorFlow als allgemeines ML-Framework, sondern den Weg von einem Modell zur robusten Edge-Inferenz mit TensorFlow Lite beziehungsweise LiteRT. Im Fokus stehen Modellkonvertierung aus Keras/SavedModel, Operator-Kompatibilität, Quantisierung, Delegate-/Beschleunigerpfade und die systematische Validierung nach der Konvertierung. Seminarrelevant wird das Thema durch die technische Leitfrage, welche Modellklassen und Operatoren sich auf Edge-Plattformen wirklich eignen, wo Konvertierungsprobleme auftreten und wie man funktionale Gleichwertigkeit, Latenz und Ressourcenbedarf sauber bewertet. Dadurch entsteht ein klar abgegrenztes Vergleichs- und Einordnungsthema statt eines breiten Framework-Tutorials.}

\TopicLeitfragen{\item Welche Modelltypen und Operatoren lassen sich für LiteRT robust konvertieren und wo liegen typische Grenzen?
\item Wie verändern Quantisierung und Hardware-Beschleunigung Latenz, Speicherbedarf und Modellgüte?
\item Wie prüft man nach der Konvertierung systematisch auf funktionale Gleichwertigkeit und numerische Abweichungen?
}
\TopicScope{\item Kein allgemeiner Einstieg in TensorFlow/Keras.
\item Kein Vollvergleich aller Edge-AI-Frameworks.
}
\TopicPractice{Ein kleines Modell nach LiteRT konvertieren, eine Float- und eine quantisierte Variante vergleichen und die Unterschiede bei Operator-Unterstützung, Laufzeit und Vorhersagen knapp dokumentieren.}

\TopicKeywords{LiteRT, TensorFlow Lite, Edge AI, Konvertierung, Quantisierung, Operatoren, Validierung}

\nocite{liteRtOverview,tfApi}
\end{topic}
